\caption{\textbf{Attention dilutes: per-sender uptake falls with the backlog $k$, and replies go to the newest messages.}
(a)~The response to a pending sender relative to a no-budget model with the same agent$\times$day propensities, vs the number $k$ of messages that reached the agent since its previous talk call (context ledger). The plot shows 16 non-holdout periods, with G51 bold and its day-bootstrap 95\% band. Every period declines, less steeply than the literal $1/k$ (dotted).
(b)~The fitted exponent $\hat\beta$ (uptake $\propto k^{-\beta}$) on all talk calls (D1, blue) is $>0$ in 16/16 periods; the pooled value is $0.66\pm0.02$ (dashed). Gray bands: generative agents with no budget whose replies are triggered (reactive timing, 8 replicates). That null alone gives D1 $\hat\beta=0.75$--$0.80$, so D1 cannot separate dilution from timing. Timer-wake batches (D2, pink: exogenous $k$) are not fooled by the null ($-0.28$ to $0.19$). D2 is clear in G51 ($0.45$ [0.41, 0.50]), G37 and G44, and it contradicts D1 in G38 and G41.
(c)~G51: the chance that a pending message becomes the reply parent, by its queue rank (1 = newest), for $k=5$--8 and 9--16, with day-bootstrap 95\% bands. Lines are the rank-blind rates.}
\label{fig:H18}
